% SPDX-License-Identifier: GPL-3.0-only
% Licensed under GNU GPL version 3; https://www.gnu.org/licenses/gpl-3.0.html . See LICENSE at repository root.
\section{The inherited LPL interface}\label{sec:foundations}

\subsection{Signatures, valuations and pure observations}
Fix a many-sorted signature $\Sig$. Its sorts and function/predicate symbols
have decidable equality; their typed argument profiles are finite. Program
variables form $\Vars=\Sort\times\N$. A model $M$ supplies nonempty carriers
$D_a$ and interpretations of the symbols. Neither carrier equality nor
predicate truth has to be decidable. The control-state type of an ASM is
finite; the sets of sorts, variables and values need not be.
\sources{\rocq{../Signature.v}{SIGNATURE},
  \rocq{../Syntax.v}{variable}, \rocq{../Semantics.v}{model}}

The code uses envelopes $E_M=\prod_{a\in\Sort}D_a$ and valuations
$\States_M=\Vars\to E_M$. Reading $x:a$ means
$\Read{s}{x}=s(x)(a)$. Term evaluation $\eval v s$ returns an envelope.
Write $s[x\leftarrow d]$ for replacing the envelope stored at $x$ by $d$.
For well-sorted equality $x=v$, only the component at $a$ is compared.
\sources{\rocq{../Semantics.v}{domain},
  \rocq{../Semantics.v}{state}, \rocq{../Semantics.v}{read},
  \rocq{../Semantics.v}{eval_term}, \rocq{../Semantics.v}{write}}

A pure formula is a first-order formula built from truth, predicates,
typed equality, negation, conjunction and universal quantification.
Other first-order connectives are abbreviations. In particular, a state
description may use derived disjunctions and existentials, but not temporal
or program modalities. For pure $A$ the two interpretations agree at the
head of any nonempty trace:
\[
 (s::\pi)\sat_M A\quad\Longleftrightarrow\quad M,s\sat_{\rm FO}A.
\]
Here $M,s\sat_{\rm FO}A$ abbreviates \code{fol_holds M A s}.
This agreement needs purity; it is not a law for arbitrary LPL formulas.
\sources{\rocq{../Syntax.v}{pureb},
  \rocq{../Semantics.v}{fol_holds}, \rocq{../Semantics.v}{pure_head}}

\subsection{Regular programs and overlap}
The instruction and continuation language has the constructors
\[
 \alpha\bnf x:=v\mid A?\mid\alpha;\beta
               \mid\alpha\cup\beta\mid\alpha^*.
\]
Composition overlaps one equal endpoint. Precisely,
$\pi\odot\eta\leadsto\theta$ means that some $\rho,s,\tau$ satisfy
\[
 \pi=\rho\cat[s],\qquad \eta=s::\tau,\qquad
 \theta=\rho\cat(s::\tau).
\]
All three traces are nonempty. For example,
$[s,s]\odot[s,t]\leadsto[s,s,t]$; the test's repeated valuation remains.
\sources{\rocq{../Syntax.v}{program},
  \rocq{../Semantics.v}{overlap_list}, \rocq{../Semantics.v}{overlap}}

For predicates $R,T$ on traces define
\[
 (R\circ T)(\theta)\iff
 \exists\pi,\eta.\ R(\pi)\land T(\eta)\land\pi\odot\eta\leadsto\theta.
\]
Then the inherited program semantics is
\begin{align*}
 \interp{x:=v}_M(\pi)&\iff
   \exists s.\ \pi=[s,s[x\leftarrow\eval v s]],\\
 \interp{A?}_M(\pi)&\iff\exists s.\ \pi=[s,s]\land[s,s]\sat_M A,\\
 \interp{\alpha;\beta}_M&=\interp\alpha_M\circ\interp\beta_M,\\
 \interp{\alpha\cup\beta}_M(\pi)&\iff
   \interp\alpha_M(\pi)\lor\interp\beta_M(\pi),\\
 \interp{\alpha^*}_M(\pi)&\iff\exists n\in\N.\ (\interp\alpha_M)^n(\pi),
\end{align*}
where $R^0(\pi)\iff\exists s.\ \pi=[s,s]$ and
$R^{n+1}=R\circ R^n$. Zero power contributes a final stutter even to
positive iterations. A single body run $[s,t]$, for instance, yields a
one-iteration star run $[s,t,t]$.
\sources{\rocq{../Semantics.v}{compose},
  \rocq{../Semantics.v}{finite_power},
  \rocq{../Semantics.v}{program_holds}}

\subsection{First, final and all-state observations}
An LPL box extends the current trace:
\[
 \pi\sat_M[\alpha]A\iff
 \forall\eta,\theta.\ \interp\alpha_M(\eta)\land
 \pi\odot\eta\leadsto\theta\ \Longrightarrow\ \theta\sat_M A
 \qquad(\pi\ne[]).
\]
A pure atom in $A$ reads the first valuation of $\theta$, not its final
valuation. In contrast, for pure $A$, the derived temporal formula
$\Alw A$ requires $A$ at every visited valuation of a nonempty trace.
The source expresses this as a \code{Forall} of singleton-trace truths;
the pure-head law gives the equivalent first-order reading.
Endpoint properties in this manual instead use the separate predicate
$\Post P\alpha$ defined in Section~\ref{sec:properties}.
\sources{\rocq{../Semantics.v}{formula_holds},
  \rocq{../Syntax.v}{Always},
  \rocq{../Semantics.v}{always_pure_forall}}
